% !TeX root = ../main.tex

\begin{abstract}

    總體基因體學的分類學分類——即將每條定序短讀序列指派至正確的微生物分類單元——是微生物群落分析的核心計算挑戰。本論文系統性評估基因體基礎模型於總體基因體分類任務的適用性，並透過資料規模、預訓練、tokenization 等多重消融實驗，分析其主要瓶頸。我們提出一種基於 Nucleotide Transformer v2（498M 參數）的 Token 層級分類方法，結合低秩適應（LoRA）微調與注意力池化分類頭，保留完整的 Token 層級序列資訊。透過系統性的資料規模實驗（500K 至 50M 條平衡讀序列），我們發現在固定 NT-v2 架構下，訓練資料量是決定分類準確度的主要因素：在屬（genus）層級 120 類分類任務中，資料量由 500K 增至 5M 可提升準確度 7.76 個百分點（55.29\% 提升至 63.05\%），進一步擴至 50M 可達 67.07\%，而架構與訓練技巧的調整最多僅帶來 $\pm$0.5 個百分點的變化。然而，將平衡資料再擴增 $5\times$ 至 250M 僅帶來 $+0.22$ 個百分點的邊際提升（達 67.29\%），顯示 NT-v2 的 6-mer 模型於 50M 之後即趨於飽和。此飽和為 tokenization 之特性而非資料之極限：以相同讀序訓練、採重疊 13-mer 的 MetaTransformer，於同一 $50\text{M}\to250\text{M}$ 區間反而由 87.5\% 持續攀升至 98.7\%，故效能上限主要由 tokenization 而非資料量決定。反向互補測試時增強（RC TTA）在所有設定下穩定提供 0.1--1.5 個百分點的準確度改善。在固定 6-mer tokenizer 下逐層掃描編碼器深度後可見，單是深度即值 15.60 個百分點（1 層至 29 層由 53.92\% 提升至 69.52\%）；因此在深度匹配且具備 50M 標註讀序的條件下，預訓練的淨貢獻趨近於零——6.45M 參數的從頭訓練模型達 69.52\%，高於 NT-v2 的 67.08\%。先前以單層模型為對照所得的 13.19 個百分點「預訓練優勢」，實際上量到的是深度。預訓練真正帶來的是資料效率：於 0.5M、1M、5M 讀序分別領先 13.4、13.1 與 8.9 個百分點，並於 50M 反轉；此交叉現象亦在獨立的 309 屬土壤目錄上重現。儘管如此，tokenization 仍是兩者中較大的因素：在 MetaTransformer 自身架構內僅更換 tokenizer 即值 40.42 個百分點，約為深度效應的 2.6 倍。以本研究之訓練配方自行訓練的 13-mer 模型中，精確 canonical 詞彙表達到 91.15\% 的屬層級準確度，而將 13-mer 雜湊至 $2^{22}$ 個桶的近似詞彙表達到 86.80\%，同時將 8~GiB 的嵌入表縮減為 2~GiB，且不同於列舉式詞彙表，可在資料平行下訓練；在寬度匹配時，雜湊的代價為 6.37 個百分點。因此 NT-v2 的非重疊 6-mer tokenizer 是本方法的主要結構性限制，而非骨幹容量或適應策略。惟此排序僅在目錄之內成立：於參考目錄之外的基因體上依 ANI 分層評估時，13-mer 模型在「已收錄物種的新菌株」上領先 32 個百分點，卻在「無任何參考基因體可及」的最遠層落後 11.3 個百分點；且封閉集準確度的保留率將兩種詞彙表分為互不重疊的兩個帶（49--54\% 對 13--22\%），深度、容量與預訓練的差異皆無法跨越。在物種（species）層級 1,535 類任務上，per-genus 50M LoRA 分類器於 oracle genus routing 下達到 29.5\% Top-1，超越 sp\_v4 平面 oracle 上限（27.4\%），驗證階層式架構方向的正確性。樣本層級評估顯示，在物種平衡的隨機分割合成樣本中，67\% 的單讀序列準確度可轉化為相對豐度估計 Pearson 相關係數 $r = 0.993$ 與豐度 $\geq$1\% 之屬的 100\% 偵測靈敏度，但此結果是在固定群落組成下的穩定性測試，於真實多樣性社群上的泛化仍待驗證。本研究為基因體基礎模型在總體基因體學分類上的應用提供了實證基礎與可擴展的實驗框架，並指出未來方向應聚焦於以重疊長 $k$-mer tokenization 預訓練的微生物基礎模型。

\end{abstract}

\begin{abstract*}

    Metagenomic taxonomic classification---assigning each sequencing read to its source organism in the microbial taxonomy---is a fundamental computational challenge in microbiome analysis. This thesis systematically evaluates the suitability of genomic foundation models (GFMs) for metagenomic read classification through controlled ablations across data scaling, pre-training, and tokenization. We propose a token-level classification approach leveraging the pre-trained Nucleotide Transformer v2 (NT-v2, 498M parameters) with Low-Rank Adaptation (LoRA) fine-tuning and an attention-pooling head that preserves the full token-level sequence information.

    Within the fixed NT-v2 framework, training data volume is the dominant factor determining classification accuracy: a 10$\times$ increase in training data yields a $+7.76$ percentage point (pp) improvement in genus-level accuracy (55.29\% to 63.05\% on 120 genera at 5M reads, scaling further to 67.07\% at 50M reads), whereas architectural and training-technique ablations produce at most $\pm$0.5~pp of variation. A further $5\times$ increase to 250M reads, however, yields only $+0.22$~pp (67.29\%): the NT-v2 6-mer model saturates beyond 50M. This saturation is a property of the tokenizer, not the data---trained on the same reads over the same $50\text{M}\to250\text{M}$ range, an overlapping 13-mer model instead scales from 87.5\% to 98.7\% genus Top-1---so the performance ceiling is set by tokenization rather than data volume. Reverse-complement test-time augmentation consistently improves accuracy by $+0.1$ to $+1.5$~pp across all experimental settings. Balanced subsampling across species reduces the number of unclassifiable taxa (F1$=$0 classes) from 9 to 0 at the 50M scale and improves macro F1 by $+5.08$~pp over the 5M baseline. Sweeping encoder depth at a fixed 6-mer tokenizer shows that depth alone is worth $+15.60$~pp (53.92\% to 69.52\% from 1 to 29 layers), so at matched depth and 50M labelled reads the net contribution of pre-training is approximately zero: a 6.45M-parameter model trained from scratch reaches 69.52\% against NT-v2's 67.08\%. An earlier single-layer ablation that appeared to isolate a $+13.19$~pp pre-training advantage was measuring depth. What pre-training does buy is data efficiency---$+13.4$, $+13.1$ and $+8.9$~pp at 0.5M, 1M and 5M reads, crossing over by 50M---and the crossover reproduces on an independent 309-genus soil catalogue.

    Tokenization is nonetheless the larger factor of the two: within MetaTransformer's own architecture, changing only the tokenizer is worth up to $+40.42$~pp, roughly $2.6\times$ what depth is worth within ours. Trained under our own recipe, an exact canonical 13-mer vocabulary reaches 91.15\% genus Top-1 and a hashed vocabulary into $2^{22}$ buckets reaches 86.80\% while replacing an 8~GiB embedding table with a 2~GiB one that, unlike an enumerated table, trains under data parallelism; hashing costs 6.37~pp at matched width. NT-v2's non-overlapping 6-mer tokenization is therefore the principal structural limitation of the proposed approach, rather than backbone capacity or adaptation strategy. That ranking, however, holds only on the catalogue: evaluated on genomes absent from the reference and stratified by ANI, a 13-mer arm leads by 32 points on a new strain of a catalogued species and trails by 11.3 where no reference genome is within reach, and retention of closed-set accuracy separates the two vocabularies into disjoint bands (49--54\% against 13--22\%) that no difference of depth, capacity or pre-training crosses. At the species level (1,535 classes), per-genus 50M LoRA classifiers under oracle-genus routing reach 29.5\% Top-1, exceeding the flat sp\_v4 oracle ceiling of 27.4\% and validating the hierarchical-classification direction. In synthetic species-balanced random-partition samples, per-read errors largely cancel at the sample level (Pearson $r = 0.993$, 100\% detection sensitivity at $\geq$1\% abundance); evaluation on real and compositionally variable communities remains future work. This work provides empirical evidence for the relative roles of data scale, pre-training, and tokenization in GFM-based metagenomic classification, and points to overlapping long-$k$ microbial pre-training as the most promising future direction.

\end{abstract*}
